% ------------------------------------------------------------------
%  Appendice C — Glossario
%  Ricerca di riferimento: ricerca 01 (termini tecnici); soprattutto i
%    capitoli in cui ogni termine è introdotto (rimando in fondo a
%    ciascuna voce)
%  Scheda nell'indice ragionato: ricerca/06_indice-ragionato.md, §6.3
%    (40–60 termini, 1.500–2.000 parole)
%
%  STATO: prima stesura (07/10/2026), approvata dall'autore (08/10/2026); 61 voci con definizione e 13
%  rinvii («vedi»); circa 2.400 parole.
%
%  REGOLA: il glossario non contiene dati, nomi o affermazioni che non
%  siano già nei capitoli; ogni definizione riprende, condensandola, la
%  formulazione del capitolo indicato nel rimando. Se un capitolo cambia
%  la definizione di un termine, va cambiata anche la voce.
%  Unico elemento nuovo: l'equivalente inglese di alcuni termini che i
%  capitoli danno solo in italiano (self-explanation, elaborative
%  interrogation, generation effect, hypercorrection effect, productive
%  failure, worked examples, dual coding theory, method of loci,
%  keyword method, cognitive load theory, expertise reversal effect,
%  self-determination theory, growth mindset, self-regulated learning,
%  implementation intentions, working memory, long-term memory,
%  encoding specificity, reconsolidation, interference, forgetting
%  curve, savings, illusion of competence, illusion of explanatory
%  depth, fluency, metacognition, effect size, meta-analysis, formative
%  assessment, feedback, active learning, learning styles, neuromyths,
%  explicit instruction, prequestions, low-stakes, blocked practice):
%  terminologia corrente della letteratura, data perché il lettore
%  possa cercare le fonti.
%  Ordine alfabetico italiano; articoli e trattini non contano.
%  Voce: impersonale o «tu» lo studente, come nei capitoli; nessuna
%  voce presenta l'autore come docente (decisione 13).
% ------------------------------------------------------------------
\chapter{Glossario}
\label{app:glossario}

%  Comandi usati solo in questa appendice
%  \voce{termine}{inglese}{definizione}{rimandi}  (inglese può essere vuoto)
%  \rinvio{termine}{voce a cui si rinvia}
%  \lettera{A}
\newcommand{\glinglese}[1]{\if\relax\detokenize{#1}\relax\else\ (\emph{#1})\fi}
%  Separatore tra le parti della voce: spazio elastico di circa
%  mezzo quadratone. Deve essere glue (non \enspace, che è un kern):
%  solo così TeX può andare a capo lì e sillabare la parola seguente.
%  Ogni voce ha anche un piccolo \emergencystretch locale: le righe
%  con termini inglesi lunghi o sigle non spezzabili (L.~170/2010)
%  altrimenti sbordano invece di allargarsi un poco.
\newcommand{\glsep}{\hskip 0.5em plus 0.25em minus 0.1em\relax}
\newcommand{\voce}[4]{%
  \par\vspace{0.4\baselineskip plus 0.15\baselineskip}%
  \begingroup\setlength{\emergencystretch}{1.5em}%
  \noindent\hangindent=1.2em\hangafter=1
  \textbf{#1}\glinglese{#2}.\glsep #3\glsep
  {\footnotesize\itshape #4.}\par\endgroup}
\newcommand{\rinvio}[2]{%
  \par\vspace{0.4\baselineskip plus 0.15\baselineskip}%
  \begingroup\setlength{\emergencystretch}{1.5em}%
  \noindent\hangindent=1.2em\hangafter=1
  \textbf{#1}:\glsep\emph{vedi} #2.\par\endgroup}
\newcommand{\lettera}[1]{%
  \par\vspace{1.1\baselineskip plus 0.4\baselineskip}%
  {\centering\color{rubrica}\scshape\spaziato[150]{\MakeLowercase{#1}}\par}%
  \nopagebreak\vspace{0.2\baselineskip}\nopagebreak}
%  Il glossario finisce dove finisce: niente spazi stirati a fondo pagina
\raggedbottom

\iniziale{Q}{uesto glossario} raccoglie i termini tecnici del libro,
con una definizione breve e il capitolo in cui se ne parla per esteso.
Tra parentesi trovi il termine inglese, quando esiste: la letteratura
scientifica sull'apprendimento è scritta quasi tutta in inglese, e con
quel nome puoi cercare le fonti. Le parole in corsivo precedute da
\emph{vedi} rimandano a un'altra voce. Le definizioni non aggiungono
nulla a ciò che dicono i capitoli: servono a ritrovare un'idea, non a
sostituirla. E, se vuoi metterle a frutto, non rileggerle: copri la
definizione e prova a dirla tu.

\small

\lettera{A}

\voce{Alternanza}{interleaving}{Mescolare, nella stessa sessione,
  esempi o problemi di tipo diverso, invece di esercitarsi su un tipo
  alla volta (lo studio \enquote{a blocchi}). Costringe a
  \emph{discriminare}, cioè a chiedersi che tipo di problema si ha
  davanti prima di chiedersi come si risolve; per questo aiuta
  soprattutto a distinguere cose che si confondono. È un principio meno
  generale del recupero e della distribuzione: non funziona per tutto}
  {Capp.~\ref{cap:alternanza}, \ref{cap:discipline}}

\voce{Apprendimento attivo}{active learning}{Nella lezione
  universitaria, tutto ciò che dà agli studenti qualcosa da fare con la
  testa durante la lezione: domande a cui tutti rispondono, problemi in
  piccoli gruppi, discussioni guidate. Si contrappone alla lezione
  tradizionale, in cui si ascolta. Conta l'attività mentale, non quella
  fisica. \emph{Vedi anche} istruzione tra pari}
  {Capp.~\ref{cap:inganni}, \ref{cap:docenti}}

\voce{Apprendimento e prestazione}{learning, performance}{La
  prestazione è ciò che si osserva durante lo studio o subito dopo:
  quanti esercizi riesci a fare, quanto ti sembra chiaro ciò che hai
  letto. L'apprendimento è il cambiamento duraturo che resta quando la
  prestazione del momento è svanita. Le due cose possono andare in
  direzioni opposte: ciò che vedi mentre studi non ti dice quanto stai
  imparando}
  {Cap.~\ref{cap:dueforze}}

\rinvio{Arte della memoria}{metodo dei luoghi; mnemotecniche}

\voce{Autodeterminazione, teoria dell'}{self-determination
  theory}{Teoria della motivazione di Edward Deci e Richard Ryan.
  Distingue la motivazione \emph{autonoma}, per ragioni che si sentono
  proprie, da quella \emph{controllata}, per ragioni che si sentono
  imposte: il voto, la paura di deludere, la pressione dei genitori. La
  prima si associa a più perseveranza; la sostengono tre bisogni:
  autonomia, competenza, relazione}
  {Cap.~\ref{cap:motivazione}}

\voce{Autoregolazione dell'apprendimento}{self-regulated
  learning}{La capacità di governare da sé il proprio studio: decidere
  che cosa studiare e quando, controllare mentre si studia se si sta
  davvero imparando, valutare poi che cosa ha funzionato e correggere il
  tiro. È la competenza che il primo anno di università richiede tutta
  insieme}
  {Capp.~\ref{cap:transizione}, \ref{cap:valutare}}

\voce{Auto-spiegazione}{self-explanation}{Spiegarsi, passaggio per
  passaggio, perché un esempio svolto procede come procede: che legge si
  usa, e perché proprio quella. Utile soprattutto nelle discipline in cui
  si impara da esempi: matematica, fisica, economia, ma anche le
  sentenze e i casi clinici}
  {Capp.~\ref{cap:capire}, \ref{cap:discipline}}

\lettera{B}

\voce{Basso rischio, prove a}{low-stakes}{Quiz, domande e prove che
  contano poco o nulla per il voto finale. Fanno recuperare senza
  l'ansia dell'esame e permettono di sbagliare in tempo, quando l'errore
  si può ancora correggere. Un peso piccolo, non nullo, aiuta a
  prenderle sul serio}
  {Capp.~\ref{cap:capire}, \ref{cap:valutare}}

\rinvio{Blocchi, studio a}{alternanza}

\voce{Brain dump}{}{Letteralmente \enquote{svuotare la testa}: la
  forma più semplice di pratica del recupero. Chiudi il libro e scrivi
  su un foglio bianco tutto ciò che ricordi; poi controlla sul testo,
  correggi, e ciò che mancava è esattamente ciò che devi ristudiare}
  {Cap.~\ref{cap:recupero}}

\lettera{C}

\voce{Carico cognitivo, teoria del}{cognitive load theory}{Teoria
  dell'istruzione di John Sweller, fondata sui limiti della memoria di
  lavoro: chi impara qualcosa di nuovo ha poche risorse, e ogni elemento
  che non serve all'apprendimento -- una slide affollata, un passaggio
  saltato perché ovvio per chi lo sa -- consuma risorse che
  servirebbero a capire. \emph{Vedi anche} esempi svolti; inversione
  dell'expertise}
  {Cap.~\ref{cap:docenti}}

\voce{Chunk, chunking}{}{Dall'inglese \emph{chunk}, \enquote{pezzo}:
  raggruppare più elementi in un'unica unità di significato, come
  sedici lettere che diventano quattro sigle note. Il limite della
  memoria di lavoro riguarda le unità di significato, e quanto grande
  sia un'unità dipende da ciò che si sa già}
  {Cap.~\ref{cap:cervello}}

\voce{Codifica, consolidamento, recupero}{encoding, consolidation,
  retrieval}{Le tre fasi di ogni ricordo. Viene \emph{codificato},
  cioè tradotto dall'esperienza in una traccia di memoria;
  \emph{consolidato}, cioè stabilizzato nelle ore e nei giorni
  successivi, in buona parte durante il sonno; e \emph{recuperato},
  cioè riportato alla coscienza quando serve}
  {Capp.~\ref{cap:cervello}, \ref{cap:corpo}}

\rinvio{Consolidamento}{codifica, consolidamento, recupero;
  riconsolidamento}

\rinvio{Criterio}{ripasso a successive riprese}

\voce{Curva dell'oblio}{forgetting curve}{La curva che descrive quanto
  si dimentica con il passare del tempo, misurata per la prima volta da
  Hermann Ebbinghaus nel 1885 e ritrovata quasi identica da Murre e Dros
  nel 2015. La dimenticanza è rapidissima nelle prime ore e nei primi
  giorni, poi rallenta. Ogni ripresa ben distanziata la rende più
  piatta}
  {Capp.~\ref{cap:vita}, \ref{cap:storia}}

\lettera{D}

\voce{Delega cognitiva}{cognitive offloading}{Usare qualcosa fuori
  dalla mente -- una lista, un promemoria, la calcolatrice,
  un'intelligenza artificiale -- per fare un lavoro che altrimenti
  farebbe la mente. Libera risorse per altro, ma ha un costo preciso: ciò
  che deleghi, non lo alleni}
  {Cap.~\ref{cap:ia}}

\voce{Difficoltà desiderabili}{desirable difficulties}{Espressione di
  Robert Bjork: le condizioni che rendono lo studio più lento e
  faticoso oggi, e il ricordo più saldo e più flessibile domani, come
  recuperare invece di rileggere, distanziare, alternare. Sono
  desiderabili solo se riesci a superarle: la fatica utile è quella di
  chi cerca e, con qualche sforzo, trova}
  {Cap.~\ref{cap:dueforze}}

\voce{Dimensione dell'effetto}{effect size}{La misura di quanto è
  grande una differenza tra due gruppi, di solito in unità di
  deviazione standard (\emph{d} di Cohen, \emph{g} di Hedges). Su questa
  scala 0 significa nessuna differenza e 0,5 è già un effetto di buona
  grandezza: lo studente medio del gruppo favorito supera circa il 69\%
  degli studenti dell'altro. Nella ricerca educativa molti interventi
  costosi producono effetti di 0,1 o 0,2}
  {Capp.~\ref{cap:recupero}, \ref{cap:metodo}}

\voce{Distribuzione, effetto della}{spacing effect}{A parità di tempo
  totale, studiare in più sessioni separate da intervalli produce un
  ricordo a lungo termine molto migliore che studiare in un'unica
  sessione. Conta distanziare; il calendario preciso conta molto meno di
  quanto si creda (intervalli crescenti e intervalli uniformi danno
  risultati simili). \emph{Vedi anche} pratica concentrata}
  {Capp.~\ref{cap:tempo}, \ref{cap:metodo}; scheda~\ref{sch:calendario}}

\voce{Doppia codifica, teoria della}{dual coding theory}{Teoria di
  Allan Paivio: la mente elabora le informazioni con due sistemi
  distinti ma collegati, uno per il linguaggio e uno per le immagini. Ciò
  che è codificato in entrambe le forme ha due vie per essere ritrovato,
  e si ricorda meglio}
  {Cap.~\ref{cap:immagini}}

\voce{DSA}{}{Disturbi specifici di apprendimento. La legge
  italiana (L.~170/2010) ne riconosce quattro: dislessia (nella lettura),
  disgrafia (nella grafia), disortografia (nella correttezza
  ortografica), discalculia (nel calcolo e nell'elaborazione dei
  numeri). Sono \enquote{specifici} perché si manifestano in presenza di
  capacità cognitive adeguate: riguardano un'abilità, non
  l'intelligenza. \emph{Vedi anche} strumenti compensativi}
  {Cap.~\ref{cap:dsa}; scheda~\ref{sch:dislessia}}

\lettera{E}

\rinvio{Effetto del test}{pratica del recupero}

\voce{Elaborazione}{elaboration}{Il lavoro mentale con cui
  un'informazione viene collegata a ciò che si sa già: chiedersi perché,
  cercare esempi, confrontare. Non conta tanto quante volte
  un'informazione viene ripetuta, ma quanto profondamente viene
  elaborata: la mente registra ciò su cui ha lavorato}
  {Capp.~\ref{cap:cervello}, \ref{cap:capire}}

\voce{Esempi svolti}{worked examples}{Problemi risolti passo per
  passo, da studiare prima di lanciarsi nei problemi da risolvere da
  soli. Servono soprattutto all'inizio di un argomento e a chi comincia;
  poi la guida va ritirata a poco a poco, da esempi completi a esempi
  incompleti, fino ai problemi autonomi}
  {Capp.~\ref{cap:discipline}, \ref{cap:docenti}}

\lettera{F}

\voce{Fallimento produttivo}{productive failure}{Affrontare un
  problema nuovo \emph{prima} che venga spiegato come si risolve, secondo
  la proposta di Manu Kapur. Può aiutare la comprensione dei concetti, a
  una condizione: che il tentativo sia seguito da una spiegazione chiara
  e completa. Il tentativo prepara il terreno; non sostituisce
  l'insegnamento}
  {Capp.~\ref{cap:capire}, \ref{cap:docenti}}

\rinvio{Familiarità}{fluidità e familiarità}

\voce{Flashcard}{}{Schede con una domanda su un lato e la risposta
  sull'altro, su carta o in un'applicazione: uno strumento eccellente
  per ciò che va ricordato con precisione (termini, formule, date,
  articoli di legge), a patto di rispondere \emph{prima} di girare la
  scheda e di riprenderle in giorni diversi}
  {Capp.~\ref{cap:recupero}, \ref{cap:tempo}}

\voce{Fluidità e familiarità}{fluency, familiarity}{Due degli indizi
  su cui la mente fonda, in automatico, la sensazione di sapere: quanto
  facilmente scorre ciò che leggi, e quanto ti sembra di averlo già
  visto. Nella vita di tutti i giorni funzionano abbastanza bene; nello
  studio si possono ottenere senza conoscenza, per esempio rileggendo.
  \emph{Vedi anche} illusione di competenza}
  {Cap.~\ref{cap:inganni}}

\voce{Forza di immagazzinamento e forza di recupero}{storage strength,
  retrieval strength}{Le due forze di ogni ricordo secondo Robert ed
  Elizabeth Bjork: quanto è radicato, cioè con quanti collegamenti a ciò
  che sai, e quanto è accessibile in questo momento. La forza di recupero
  oscilla; quella di immagazzinamento, secondo la teoria, può soltanto
  crescere, e cresce tanto di più quanto più la forza di recupero è
  bassa nel momento in cui richiami qualcosa. Mentre studi vedi
  soprattutto la seconda forza, non la prima}
  {Capp.~\ref{cap:dueforze}, \ref{cap:inganni}}

\lettera{G}

\voce{Generazione, effetto di}{generation effect}{Ciò che si produce
  da sé si ricorda meglio di ciò che si legge: chi deve completare
  \enquote{caldo -- f\dots} ricorda \emph{freddo} meglio di chi legge la
  coppia intera. Lo sforzo di produrre una risposta, anche facile, lascia
  una traccia più profonda. \emph{Vedi anche} pre-domande}
  {Cap.~\ref{cap:capire}}

\rinvio{Grinta}{mentalità di crescita e grinta}

\voce{Gruppo di controllo}{control group}{In un esperimento, il gruppo
  che non riceve l'intervento (o ne riceve uno diverso) e con cui si
  confronta chi lo riceve: senza confronto non si può sapere se un
  miglioramento si deve all'intervento. Quando chi entra in un gruppo o
  nell'altro è deciso a caso, lo studio si dice \emph{randomizzato}}
  {Capp.~\ref{cap:sapere}, \ref{cap:alternanza}}

\lettera{I}

\voce{Illusione della profondità esplicativa}{illusion of explanatory
  depth}{Crediamo di capire le cose molto meglio di quanto le capiamo,
  e ce ne accorgiamo soltanto quando dobbiamo spiegarle. Il rimedio è
  appunto provare a spiegare}
  {Capp.~\ref{cap:inganni}, \ref{cap:insegnare}}

\voce{Illusione di competenza}{illusion of competence}{Giudicare ciò
  che sapremo fare \emph{senza} la risposta in base a come ci sentiamo
  \emph{con} la risposta davanti. Nasce dall'alta forza di recupero di
  ciò che si è appena letto. Per sapere che cosa si sa bisogna
  verificarlo: a libro chiuso, con la sola domanda, a distanza di tempo}
  {Cap.~\ref{cap:inganni}}

\voce{Insegnare per imparare}{learning by teaching}{Prepararsi a
  spiegare, e spiegare davvero, fa imparare chi spiega: chi si aspetta di
  dover insegnare legge cercando la struttura, e chi insegna a memoria,
  senza appunti, recupera ad alta voce. Insegnare leggendo da un testo,
  invece, non aiuta. È il \emph{docendo discimus} degli antichi}
  {Cap.~\ref{cap:insegnare}}

\voce{Intenzioni di implementazione}{implementation intentions}{Decidere
  in anticipo \emph{quando} e \emph{dove} si farà \emph{che cosa},
  nella forma \enquote{quando accade X, faccio Y}: \enquote{il martedì,
  appena finita la lezione, vado in biblioteca e faccio il primo
  recupero}. È il rimedio più semplice contro il rinvio}
  {Capp.~\ref{cap:transizione}, \ref{cap:dsa}}

\voce{Interferenza}{interference}{Ricordi simili si ostacolano a
  vicenda: il nuovo PIN con il vecchio, lo spagnolo con il francese, gli
  articoli del codice studiati nella stessa settimana. Più le cose sono
  simili e meno sono distinte nella mente, più si disturbano. È una delle
  cause principali dell'oblio}
  {Cap.~\ref{cap:cervello}}

\voce{Interrogazione elaborativa}{elaborative interrogation}{Fermarsi,
  mentre si studia, e chiedersi \emph{perché}: perché questa cosa è
  vera, perché accade così e non altrimenti, come si collega a ciò che
  so già. Il nome è pesante; l'abitudine è quella dei bambini di tre
  anni}
  {Capp.~\ref{cap:capire}, \ref{cap:metodo}}

\voce{Inversione dell'expertise}{expertise reversal effect}{Ciò che
  aiuta il principiante -- la guida dettagliata, l'esempio svolto, la
  ridondanza -- può ostacolare chi è già esperto, per il quale diventa
  un peso inutile. La didattica migliore cambia con il livello di chi
  impara}
  {Capp.~\ref{cap:docenti}, \ref{cap:miti}}

\voce{Ipercorrezione, effetto di}{hypercorrection effect}{Gli errori
  commessi con grande sicurezza, quando vengono corretti, si correggono
  in modo particolarmente stabile: la sorpresa di aver sbagliato ciò di
  cui si era certi fissa la risposta giusta}
  {Cap.~\ref{cap:capire}}

\voce{Ippocampo}{hippocampus}{Struttura dalla forma ricurva, nel lobo
  temporale, da cui dipende la formazione di nuovi ricordi. Durante il
  sonno profondo riattiva le tracce del giorno e le trasferisce a poco a
  poco alla corteccia. Il paziente H.M., che lo aveva perduto, non
  riusciva più a formare nuovi ricordi dichiarativi}
  {Capp.~\ref{cap:cervello}, \ref{cap:corpo}}

\voce{Istruzione esplicita}{explicit instruction}{Insegnamento per
  passi chiari, con molta esercitazione e correzione immediata. Combinata
  con l'insegnamento di strategie, è tra gli interventi più efficaci per
  chi ha un disturbo dell'apprendimento; non è in contrasto con il
  fallimento produttivo, che cambia l'ordine ma non toglie la
  spiegazione}
  {Capp.~\ref{cap:dsa}, \ref{cap:docenti}}

\voce{Istruzione tra pari}{peer instruction}{Il metodo reso celebre da
  Eric Mazur: una domanda concettuale a cui ciascuno risponde da solo;
  due minuti per convincere il vicino della propria risposta; una nuova
  votazione e la spiegazione. Fa recuperare, e fa imparare chi spiega}
  {Capp.~\ref{cap:insegnare}, \ref{cap:docenti}}

\lettera{M}

\voce{Memoria a lungo termine}{long-term memory}{Il magazzino di tutto
  ciò che sai, di cui nessuno ha trovato il fondo. Vi convivono almeno
  due sistemi: la memoria \emph{dichiarativa}, di ciò che si può dire,
  che comprende la memoria \emph{semantica} (fatti e concetti) e quella
  \emph{episodica} (episodi vissuti); e la memoria \emph{procedurale}, di
  ciò che si sa fare senza bisogno di dirlo}
  {Cap.~\ref{cap:cervello}}

\voce{Memoria di lavoro}{working memory}{Lo spazio mentale in cui
  teniamo presenti le cose mentre ci ragioniamo sopra. Contiene
  all'incirca quattro elementi alla volta e li trattiene per pochi
  secondi; ciò che si sa già, ben organizzato, le permette di trattare
  molte cose come un'unica unità. \emph{Vedi anche} chunk; carico
  cognitivo}
  {Capp.~\ref{cap:sapere}, \ref{cap:cervello}, \ref{cap:attenzione}}

\voce{Mentalità di crescita e grinta}{growth mindset, grit}{Due idee di
  grande successo. La prima, resa celebre da Carol Dweck: chi crede che
  l'intelligenza si possa sviluppare affronterebbe meglio le difficoltà.
  La seconda, di Angela Duckworth: la combinazione di passione e
  perseveranza conterebbe più del talento. Contengono qualcosa di vero,
  ma le prove raccolte invitano a ridimensionarle molto}
  {Cap.~\ref{cap:motivazione}}

\voce{Meta-analisi}{meta-analysis}{Una ricerca che raccoglie tutti gli
  studi disponibili su un tema e ne calcola l'effetto medio. È lo
  strumento con cui questo libro misura la solidità di una strategia; ma
  anche le stime più famose, rifatte con più cura, possono
  ridimensionarsi. \emph{Vedi anche} dimensione dell'effetto}
  {Capp.~\ref{cap:recupero}, \ref{cap:valutare}}

\voce{Metacognizione}{metacognition}{Il giudizio su come funziona la
  propria mente: su ciò che si sa e ciò che non si sa, e su quali
  strategie funzionano. Fidarsi della sensazione di sapere è un errore di
  metacognizione, il primo dei tre errori da cui parte il libro}
  {Capp.~\ref{cap:vita}, \ref{cap:inganni}}

\voce{Metodo dei luoghi}{method of loci}{La mnemotecnica che la
  tradizione attribuisce a Simonide, detta anche \emph{palazzo della
  memoria}: si sceglie un percorso di luoghi in un ordine fisso, vi si
  collocano immagini delle cose da ricordare, e l'ordine dei luoghi
  conserva l'ordine delle cose. È il cuore dell'antica arte della
  memoria; serve per le sequenze ordinate}
  {Capp.~\ref{cap:storia}, \ref{cap:immagini}}

\voce{Mnemotecniche}{mnemonics}{Tecniche che creano un aggancio
  artificiale per un'informazione che, di per sé, non ne ha: acronimi,
  rime, il metodo dei luoghi, il metodo della \emph{parola chiave} per i
  vocaboli stranieri (il tedesco \emph{Tisch}, \enquote{tavolo}, e una
  tisana rovesciata su un tavolo). Aiutano a ritrovare un'informazione,
  non a capirla}
  {Cap.~\ref{cap:immagini}}

\rinvio{Motivazione autonoma e controllata}{autodeterminazione, teoria
  dell'}

\voce{Multitasking}{}{Fare più cose insieme. La mente può fare insieme
  una cosa impegnativa e una automatica, ma non due cose impegnative:
  quando sembra che lo faccia, in realtà alterna, e ogni passaggio costa
  tempo, errori e profondità. Costa a tutti, \enquote{nativi digitali}
  compresi}
  {Capp.~\ref{cap:attenzione}, \ref{cap:miti}}

\lettera{N}

\voce{Neuromiti}{neuromyths}{Credenze false sul cervello e
  sull'apprendimento, diffuse anche tra gli insegnanti: gli stili di
  apprendimento, il cervello usato al dieci per cento, la musica di
  Mozart che rende più intelligenti, il \emph{brain training}. Chi sa
  di più sul cervello non ne è immune: il rimedio non è sapere meno di
  neuroscienze, ma leggerle meglio}
  {Capp.~\ref{cap:docenti}, \ref{cap:miti}, \ref{cap:formatori}}

\lettera{P}

\rinvio{Palazzo della memoria}{metodo dei luoghi}

\voce{Permastore}{}{Il nome, \enquote{deposito permanente}, che Harry
  Bahrick ha dato alla parte della memoria che sopravvive quasi intatta
  per decenni. Quanto vi entra non dipende dal caso, ma da quanto bene si
  è imparato all'inizio}
  {Cap.~\ref{cap:vita}}

\voce{Pratica concentrata}{massed practice}{Studiare tutto in
  un'unica sessione invece che in più sessioni distanziate; nella forma
  estrema, quella della vigilia dell'esame, gli studenti di lingua
  inglese la chiamano \emph{cramming}, \enquote{stipare}. Dà una buona
  prestazione il giorno dopo e un ricordo scarso nel tempo.
  \emph{Vedi anche} distribuzione}
  {Cap.~\ref{cap:tempo}}

\voce{Pratica del recupero}{retrieval practice, testing effect}{Richiamare
  dalla memoria ciò che si è studiato -- a libro chiuso, con domande,
  flashcard, esercizi senza soluzione -- invece di rileggerlo. Ciò che
  rafforza un ricordo non è il tempo passato davanti al libro, ma lo
  sforzo di richiamarlo: il test non serve soltanto a misurare quanto si
  è imparato, ma a imparare (\enquote{effetto del test}). È il modo più
  potente di studiare che si conosca}
  {Capp.~\ref{cap:recupero}, \ref{cap:metodo}; scheda~\ref{sch:metodo}}

\rinvio{Pratica distribuita}{distribuzione, effetto della}

\voce{Pre-domande}{prequestions}{Domande a cui provi a rispondere
  \emph{prima} di leggere un capitolo o di seguire una lezione, anche se
  non conosci ancora l'argomento e sbaglierai quasi certamente.
  Migliorano l'apprendimento delle informazioni su cui vertevano, non
  del resto del materiale}
  {Capp.~\ref{cap:capire}, \ref{cap:semestre}}

\rinvio{Prestazione}{apprendimento e prestazione}

\lettera{Q}

\voce{Quaderno degli errori}{}{Il quaderno in cui annoti ogni errore
  commesso negli esercizi, insieme alla sua specie: di
  \emph{conoscenza} (non sapevi la definizione o la formula), di
  \emph{riconoscimento} (non hai capito quale strumento usare), di
  \emph{esecuzione} (un segno sbagliato, un passaggio saltato). Ogni
  specie ha il suo rimedio}
  {Capp.~\ref{cap:semestre}, \ref{cap:discipline}}

\lettera{R}

\rinvio{Randomizzato, studio}{gruppo di controllo}

\rinvio{Recupero}{pratica del recupero; codifica, consolidamento,
  recupero}

\voce{Riconsolidamento}{reconsolidation}{Un ricordo consolidato,
  quando viene richiamato, torna per qualche tempo fragile e va
  stabilizzato di nuovo. Ogni volta che richiami un ricordo, dunque, lo
  riscrivi: se lo richiami bene, e lo correggi quando sbagli, esce
  rafforzato e aggiornato; se lo richiami male e nessuno ti corregge,
  puoi consolidare anche un errore}
  {Cap.~\ref{cap:cervello}}

\voce{Ripasso a successive riprese}{successive relearning}{La
  combinazione di recupero e distribuzione proposta da Katherine Rawson
  e John Dunlosky. In ogni sessione ti interroghi a libro chiuso e
  rimetti in fondo al mazzo le domande sbagliate, finché ciascuna non ha
  avuto almeno una risposta giusta (il \enquote{criterio}); poi ripeti
  la stessa procedura in più sessioni distanziate, di solito da tre a
  cinque}
  {Capp.~\ref{cap:tempo}, \ref{cap:metodo}}

\voce{Riscontro}{feedback}{L'informazione su com'è andata una risposta
  o una prova. Rende di più quando dice perché e come migliorare, invece
  di limitarsi a \enquote{giusto} o \enquote{sbagliato}; e serve quando
  arriva in tempo per correggersi. Il riscontro che arriva tardi, come il
  voto settimane dopo lo studio, tiene in vita le strategie ingannevoli}
  {Capp.~\ref{cap:inganni}, \ref{cap:valutare}}

\voce{Risparmio}{savings}{La misura inventata da Ebbinghaus: quanto
  lavoro in meno serve per reimparare ciò che si era imparato. Mostra
  che un ricordo apparentemente perduto, di cui non si sa ripetere
  nulla, si reimpara più in fretta di uno mai posseduto: non è sparito
  del tutto}
  {Capp.~\ref{cap:storia}, \ref{cap:dueforze}}

\lettera{S}

\voce{Schema}{schema}{Una struttura di conoscenze già organizzate a
  partire dalla quale si comprende e si ricorda. Frederic Bartlett ne
  concluse che ricordare non è riprodurre, ma ricostruire; e si ricorda
  meglio ciò che si riesce a legare a ciò che si conosce}
  {Capp.~\ref{cap:storia}, \ref{cap:capire}}

\voce{Sessione settimanale di ripresa}{}{Il cuore del metodo del
  libro: per ogni insegnamento, un appuntamento fisso alla settimana, da
  mezz'ora a un'ora, scritto in agenda come una lezione. Si risponde a
  libro chiuso, fino al criterio, alle domande di tutte le settimane
  precedenti; quelle che si sanno bene tornano più di rado, quelle
  sbagliate subito}
  {Cap.~\ref{cap:metodo}; schede~\ref{sch:metodo} e~\ref{sch:calendario}}

\voce{Specificità della codifica}{encoding specificity}{Principio di
  Endel Tulving: un ricordo si ritrova più facilmente quando le
  condizioni del recupero somigliano a quelle in cui è stato formato.
  Nello studio reale l'effetto è piccolo, e la conseguenza utile è il
  contrario di ciò che sembra: ritrovare le stesse cose in modi, momenti
  e contesti diversi rende il ricordo più libero}
  {Cap.~\ref{cap:cervello}}

\voce{Stili di apprendimento}{learning styles}{L'idea che ognuno sia
  \enquote{visivo}, \enquote{uditivo} o \enquote{cinestetico} e impari
  meglio se l'insegnamento corrisponde al suo stile. È un mito: quasi
  nessuno studio ha usato il disegno sperimentale necessario a
  dimostrarlo, e quelli che lo hanno usato danno risultati contrari.
  Le preferenze esistono; non dicono come si impara meglio}
  {Capp.~\ref{cap:dsa}, \ref{cap:miti}}

\voce{Strumenti compensativi e misure dispensative}{}{Per gli
  studenti con DSA, i primi compensano una difficoltà strumentale
  (sintesi vocale, testi digitali, computer con correttore,
  calcolatrice, mappe); le seconde dispensano da una prestazione o ne
  cambiano le condizioni (tempo aggiuntivo, scelta tra scritto e orale,
  valutazione del contenuto più che della forma). Il criterio che il
  libro propone: \emph{compensare l'accesso, non il pensiero}}
  {Cap.~\ref{cap:dsa}; scheda~\ref{sch:dislessia}}

\lettera{V}

\voce{Valutazione formativa}{formative assessment}{Usare le prove
  durante il corso per capire a che punto sono gli studenti e adattare
  l'insegnamento, secondo la proposta di Paul Black e Dylan Wiliam.
  L'idea resta buona; ma le prime stime del suo effetto, molto alte, si
  sono ridimensionate quando i conti sono stati rifatti}
  {Cap.~\ref{cap:valutare}}

\normalsize

\finecapitolo
